%%%%%%%%%%%%%%%%%%%%%%%%%%%%%%%%%%%%%%%%%%%%%%%%%%%%%%%%%%%%%%%%%%%%%%
% Lớp 10 | Xác suất | Xác suất của biến cố
% Dạng 0D0?2-3: Liên quan việc sắp xếp chỗ
% Số câu: 21  (sinh tự động bởi tach_cau.py; từ nay thầy sửa trực tiếp file này)
%%%%%%%%%%%%%%%%%%%%%%%%%%%%%%%%%%%%%%%%%%%%%%%%%%%%%%%%%%%%%%%%%%%%%%

%%%%%%%%%%%%%%%% TRẮC NGHIỆM 4 PHƯƠNG ÁN %%%%%%%%%%%%%%%%

% nguồn: MaHoaNH_TracNghiem_Lop10
\begin{ex}%[0D0H2-3]%[ma:5b7b5f1a]
	Gieo  hai con xúc xắc cân đối và đồng chất. Tính xác suất để hiệu số chấm xuất hiện trên hai con xúc xắc đó bằng $2$.
	\choice 
		{$\dfrac{1}{3}$} 
		 {$\dfrac{1}{7}$}
		 {\True $\dfrac{2}{9}$}
		{$\dfrac{3}{8}$} 
	\loigiai{ 
		Tập hợp  $\Omega$ có $36$ phần tử.\\
		Gọi $D$  là biến cố \lq\lq hiệu số chấm xuất hiện trên hai con xúc xắc bằng $2$\rq\rq.\\
		Các kết quả thuận lợi cho $D$ là: $(1;3), (3;1), (2;4), (4;2), (4;6), (6;4),  (3;5), (5;3)$.\\
		Tập hợp $D$  có $8$ phần tử.\\
		Vậy xác suất biến cố $D$ là $P(D) = \dfrac{n(D)}{n(\Omega)} = \dfrac{8}{36} =  \dfrac{2}{9}$.
	 }
\end{ex}

% nguồn: MaHoaNH_TracNghiem_Lop10
\begin{ex}%[0D0H2-3]%[ma:fbefdaab]
	Một nhóm gồm $7$ học sinh nam và $4$ học sinh nữ xếp thành một hàng ngang để tham gia một trò chơi. Tính xác suất để khi  xếp $2$ học sinh nữ bất kì không đứng cạnh nhau.
	\choice 
		{$\dfrac{4}{33}$} 
		{$\dfrac{3}{7}$} 
		 {$\dfrac{8}{11}$}
		{\True $\dfrac{7}{33}$} 
	\loigiai{ 
		Số phần tử của không gian mẫu là số cách xếp bất kì $11$ học sinh trên thành một hàng ngang nên $n\left(\Omega\right)=11!$.\\ 
		Gọi $A$ là biến cố \lq\lq Khi xếp hàng $2$ học sinh  nữ bất kì không đứng cạnh nhau\rq\rq.\\
		Xếp $7$ học sinh nam trước có  $7!$ cách.\\
		Khi đó $7$ học sinh này xếp hàng tạo $8$ khoảng trống. 
		Xếp $4$ học sinh nữ, mỗi bạn vào một khoảng trống sao cho không có $2$ bạn nào cùng khoảng trống có $\mathrm{A}_8^4$ cách.\\ 
		Suy  ra $n(A)=7!\cdot \mathrm{A}_8^4$.\\ $\mathrm{P}(A)=\dfrac{n(A)}{n\left(\Omega\right)}
		 =\dfrac{7!\cdot \mathrm{A}_8^4}{11!}=\dfrac{7}{33}$.
	 }
\end{ex}

% nguồn: MaHoaNH_TracNghiem_Lop10
\begin{ex}%[0D0H2-3]%[ma:1885cbc4]
	Gieo hai con xúc xắc.  Xác suất của biến cố \lq\lq Tổng số chấm trên hai mặt chia hết cho $3$\rq\rq \, là
	\choice 
		 {$\dfrac{11}{36}$}
		 {$\dfrac{13}{36}$}
		{\True $\dfrac{1}{3}$} 
		 {$\dfrac{1}{4}$}
	 \loigiai{
		 Không gian mẫu là $\Omega = \{(i,j) \mid i,j \in \{1,2,3,4,5,6\}\}$, có
		\[ 
		n(\Omega)  = 6 \cdot 6 = 36.
		\] 
		Xét biến cố $A$: \lq\lq Tổng số chấm trên hai mặt  chia hết cho $3$\rq\rq.\\
		$A=\left\{(1;2), (2;1), (2;4),  (4;2), (3;6), (6;3), (4;5), (5;4),(1;5), (5;1),(3,\;3), (6;6)\right\}\Rightarrow n(A)=12$.
		\[\Rightarrow  \mathrm{P}(A) = \dfrac{12}{36} = \dfrac{1}{3}.
		\] 
	 }
\end{ex}

% nguồn: MaHoaNH_TracNghiem_Lop10
\begin{ex}%[0D0H2-3]%[ma:62118db6]
	Xếp ngẫu nhiên $2$ bạn nam và $2$ bạn nữ ngồi vào $4$ ghế kê theo hàng ngang. Tính xác suất sao cho nam, nữ ngồi xen kẽ  nhau.
	 \choice
		 {$\dfrac{1}{2}$}
		{$\dfrac{2}{3}$} 
		 {\True $\dfrac{2}{3}$}
		 {$\dfrac{1}{4}$}
	\loigiai{ 
		Số phần tử của không  gian mẫu là $n(\Omega)=4!=24$.\\
		Gọi $A$ là biến  cố: \lq \lq Nam, nữ ngồi xen kẽ nhau\rq \rq.\\
		Giả sử đánh  số ghế từ $1$ đến $4$.\\
		\textbf{Trường hợp 1:} Nam ngồi ghế số $1$ và $3$, suy ra có $2!$ cách  xếp. \\
		Nữ ngồi ghế  số $2$ và $4$, suy ra có $2!$ cách xếp.\\
		Vậy trường  hợp 1 có $2! \cdot 2!=4$ cách xếp.\\
		\textbf{Trường hợp 2:} Nữ ngồi ghế số $1$ và $3$, suy ra có $2!$ cách  xếp. \\
		Nam ngồi ghế số $2$ và $4$, suy ra có $2!$ cách  xếp.\\
		Vậy trường hợp  $2$ có $2! \cdot 2!=4$ cách xếp.\\
	Vậy $n(A)=4+4=8$. Do  đó $ \mathrm{P}(A)=\dfrac{8}{24}=\dfrac{1}{3}$.}
\end{ex}

% nguồn: MaHoaNH_TracNghiem_Lop10
\begin{ex}%[0D0H2-3]%[ma:27606795]
	Một bộ bài tú lơ khơ gồm  $52$ quân. Lấy ngẫu nhiên $3$ quân bài. Xác suất để lấy được ít nhất $2$ quân $4$ là
	 \choice
		 {$\dfrac{288}{5525}$}
		{$\dfrac{76}{5525}$} 
		 {\True $\dfrac{73}{5525}$}
		{$\dfrac{12}{221}$} 
	 \loigiai{
		Số phần tử của không gian  mẫu $n(\Omega) = \mathrm{C}^3_{52}$.\\
		Xét biến cố $A$: \lq\lq lấy được ít nhất  $2$ quân $4$\rq\rq.\\
		\textbf{Trường hợp 1:} Lấy đúng $2$ quân $4$, còn lại $1$ lá không phải $4$ có $\mathrm{C}_4^2  \cdot \mathrm{C}^1_{48} $ cách.\\
		\textbf{Trường hợp 2:} Lấy đúng  $3$ quân $4$ có $ \mathrm{C}^3_4$ cách.\\
		Vậy 
		\[ 
		n(A) = \mathrm{C}_4^2  \cdot \mathrm{C}^1_{48}+ \mathrm{C}^3_4.
		\] 
		\[ 
		\mathrm{P}(A) =  \dfrac{n(A)}{n(\Omega)} = \dfrac{\mathrm{C}_4^2 \cdot \mathrm{C}^1_{48}+ \mathrm{C}^3_4}{\mathrm{C}^3_{52}} = \dfrac{73}{5525}.
		\] 
	 }
\end{ex}

% nguồn: MaHoaNH_TracNghiem_Lop10
\begin{ex}%[0D0H2-3]%[ma:00a79d36]
	Tung một đồng xu cân đối và đồng chất liên tiếp cho đến khi xuất hiện mặt sấp hoặc cả năm  lần đều là ngửa thì dừng lại. Tính xác suất để số lần tung không vượt quá bốn.
	 \choice
		 {$\dfrac{2}{5}$}
		 {$\dfrac{3}{5}$}
		 {\True $\dfrac{2}{3}$}
		{$\dfrac{1}{3}$} 
	 \loigiai{
		Không gian mẫu gồm các kết  quả dừng lại khi gặp mặt sấp hoặc $5$ lần ngửa
		\[ 
		\Omega = \{\mathrm{S}, \mathrm{NS}, \mathrm{NNS}, \mathrm{NNNS},  \mathrm{NNNNS}, \mathrm{NNNNN}\}\Rightarrow n(\Omega) = 6.
		 \]
		Biến cố $A$:  \lq\lq Số lần tung không vượt quá $4$\rq\rq.
		 \[
		 A = \{\mathrm{S}, \mathrm{NS}, \mathrm{NNS}, \mathrm{NNNS}\}\Rightarrow n(A) = 4.
		 \]
		Vậy xác suất của biến cố $A$ là 
		\[ 
		\mathrm{P}(A) = \dfrac{4}{6} = \dfrac{2}{3}. 
		\] 
	 }
\end{ex}

% nguồn: MaHoaNH_TracNghiem_Lop10
\begin{ex}%[0D0H2-3]%[ma:e312dea2]
	Một bộ bài tú lơ khơ gồm $52$ quân. Lấy ngẫu nhiên $3$ quân bài. Xác suất để $3$ quân bài rút ra có một  con $2$, một con $4$ và một con K là
	 \choice
		 {$\dfrac{3}{5525}$}
		{\True  $\dfrac{16}{5525}$}
		 {$\dfrac{7}{11050}$}
		{$\dfrac{8}{5525}$} 
	\loigiai{ 
		Số phần tử của không  gian mẫu là số cách chọn $3$ lá bài trong $52$ lá
		\[ 
		 n(\Omega) = \mathrm{C}^3_{52}.
		 \]
		Biến cố $A$: \lq\lq $3$ quân bài  rút ra gồm một con $2$, một con $4$ và một con K\rq\rq.
		 \[
		n(A) = \mathrm{C}^1_4 \cdot \mathrm{C}^1_4 \cdot \mathrm{C}^1_4 = 4 \cdot  4 \cdot 4 = 64.
		\] 
		Vậy xác suất của biến cố  $A$ là
		 \[
		\mathrm{P}(A)  = \dfrac{64}{\mathrm{C}^3_{52}} = \dfrac{16}{5525}.
		\] 
	 }
\end{ex}

% nguồn: MaHoaNH_TracNghiem_Lop10
\begin{ex}%[0D0V2-3]%[ma:ae2a3ca6]
	Có $7$ chiếc ghế được kê thành một hàng ngang. Xếp ngẫu nhiên $7$ học sinh, gồm $3$ học sinh lớp  $A$, $2$ học sinh lớp $B$ và $2$ học sinh lớp $C$, ngồi vào hàng ghế đó, sao cho mỗi ghế có đúng một học sinh. Xác suất để $2$ học sinh lớp $C$ không ngồi cạnh nhau và cũng không ngồi cạnh học sinh lớp $A$ bằng
	 \choice
		{$\dfrac{\left(2\cdot2\cdot3\right)!}{7!}$} 
		 {$\dfrac{2!\cdot2!}{7!}$}
		{$\dfrac{1}{70}$} 
		{\True $\dfrac{1}{105}$} 
	\loigiai{ 
		Xếp tất cả $7$ học sinh vào $7$ ghế theo một hàng ngang, số  phần tử không gian mẫu $n\left(\Omega\right)=7!$.\\
		Gọi $A$ là biến cố \lq\lq $2$ học sinh lớp $C$ không ngồi  cạnh nhau và cũng không ngồi cạnh học sinh lớp $A$\rq\rq.\\
		Xét các trường hợp xảy ra của biến  cố $A$.
		\begin{itemize}
			\item {\bf Trường hợp 1:} Xếp học $1$ sinh lớp $C$ ở ghế thứ nhất  như thế ghế thứ hai là học sinh lớp $B$ ghế thứ $3$ là học sinh lớp $C$ ghế thứ $4$ là học sinh lớp $B$ các ghế còn lại là học sinh lớp $A$.
			\begin{center}
				\begin{tabular}{|c|c|c|c|c|c|c|c|}
					 \hline
					Ghế số &  1 & 2 & 3 & 4 & 5 & 6 & 7 \\
					\hline 
					Học  sinh lớp & C & B & C & B & A & A & A \\
					\hline 
				\end{tabular}
			\end{center}
			Vậy có  $2\cdot1\cdot2\cdot1\cdot3!=12$ (cách).
			\item {\bf Trường hợp 2:} Xếp học $1$ sinh lớp $C$ ở ghế thứ $7$ như  thế ghế thứ $6$ là học sinh lớp $B$ ghế thứ $5$ là học sinh lớp $C$ ghế thứ $4$ là học sinh lớp $B$ các ghế còn lại là học sinh lớp $A$.
			\begin{center}
				\begin{tabular}{|c|c|c|c|c|c|c|c|}
					 \hline
					Ghế số & 1 & 2 & 3 & 4 & 5 & 6  & 7 \\
					\hline 
					Học sinh lớp & A & A & A & B & C  & B & C \\
					 \hline
				\end{tabular}
			\end{center}
			Vậy có  $ 2\cdot1\cdot2\cdot1\cdot3!=12$ (cách).
			\item {\bf Trường hợp 3:} Xếp học sinh lớp $C$ lần lượt tại vị trí $1$ và $7$, học sinh lớp  $B$ lần lượt tại vị trí $2$ và $6$ khi đó $3$ học sinh lớp $A$ xếp vào các vị trí còn lại.
			\begin{center}
				\begin{tabular}{|c|c|c|c|c|c|c|c|}
					 \hline
					Ghế số & 1 & 2  & 3 & 4 & 5 & 6 & 7 \\
					\hline 
					 Học sinh lớp & C & B & A & A & A & B & C \\
					 \hline
				\end{tabular}
			\end{center}
			 Vậy có $ 2!\cdot2!\cdot3!=24$ (cách).
		\end{itemize}
		Vậy số phần tử biến cố $A$ là $n(A)=12+12+24=48$  (cách).\\
		Xác suất  biến cố $D$ là $\mathrm{P}(A)=\dfrac{n(A)}{n\left(\Omega\right)}=\dfrac{48}{7!}=\dfrac{1}{105}$.
	} 
\end{ex}

% nguồn: MaHoaNH_TracNghiem_Lop10
\begin{ex}%[0D0C2-3]%[ma:d49436bc]
	Gọi $S$ là tập hợp tất cả các số tự nhiên gồm $4$ chữ số đôi một khác nhau được chọn từ các chữ số $2,3,4,5,6,7,8,9$. Xác định  số phần tử của $S$. Lấy ngẫu nhiên một số tự nhiên từ $S$, tính xác suất để số được chọn là số chia hết cho $11$ và tổng $4$ chữ số của nó cũng chia hết cho $11$.
	\choice 
		 {$\dfrac{6}{35}$}
		 {$\dfrac{1}{70}$}
		 {$\dfrac{1}{140}$}
		{\True $\dfrac{1}{35}$} 
	 \loigiai{
		Số phần  tử của $S$ là $\mathrm{A}_8^4$.\\
		Ký hiệu số có $4$ chữ số là $\overline{abcd} =  1000a + 100b + 10c + d = 100\overline{ab} + \overline{cd}$.\\
		Ta viết lại: 
		 \[
		\overline{abcd} = 11\cdot (9ab + a + c) + (b + d - a  - c)
		\] 
		Do đó, để số  chia hết cho $11$ thì $(b + d - a - c)$ phải chia hết cho $11$.\\
		Đồng thời, giả thiết yêu cầu tổng $a + b + c + d$ cũng  chia hết cho $11$.\\
		Gọi $T = a +  b + c + d$ thì điều kiện trở thành
		\[ 
		\heva{& (b + d - a - c) \equiv 0 \pmod{11} \\&(a + b + c + d)  \equiv 0 \pmod{11}}
		 \Rightarrow
		\heva{& b +  d \equiv a + c \pmod{11} \\&b + d + a + c \equiv 0 \pmod{11}}
		\Rightarrow \heva{&(b + d) \vdots 11\\& (a +  c)\vdots 11.}
		\] 
		Do $3\leq b 
		+d\le 17$,  $3\le a+c\le 17$ nên $b+d=a+c=11$.\\
		Ta xét các cặp $(a + c, b + d)$ sao  cho tổng bằng $11$ là\\
		$(2,9), (3,8), (4,7),  (5,6)\Rightarrow 4$ cặp.\\
		Số lượng các số thỏa  mãn yêu cầu là
		$A_4^2 \cdot 2!  \cdot 2! $
		$\Rightarrow \mathrm{P}(A) =  \dfrac{\mathrm{A}_4^2 \cdot 2! \cdot 2! }{\mathrm{A}_8^4} = \dfrac{1}{35}$.
	} 
\end{ex}


%%%%%%%%%%%%%%%% ĐÚNG/SAI %%%%%%%%%%%%%%%%

% nguồn: MaHoaNH_DungSai_Lop10
\begin{ex}%[0D0V2-3]%[ma:ac14bdb3]
	Cho $12$ học sinh ($9$ nam, $3$ nữ).  Xếp ngẫu nhiên thành hàng ngang. Xác suất để 3 nữ cạnh nhau là
	 \choiceTF
		 {$\dfrac{1}{220}$}
		 {$\dfrac{1}{132}$}
		{\True  $\dfrac{1}{22}$}
		{$\dfrac{1}{111}$} 
	\loigiai{ 
		Số cách xếp ngẫu nhiên $12$ học sinh thành hàng  ngang là $12!$.\\
		Để xếp  $12$ học sinh thành hàng ngang sao cho $3$ học sinh nữ cạnh nhau ta thực hiện như sau:
		\begin{itemize}
			 \item Xếp 3 nữ cạnh nhau được nhóm $\mathrm{N}$, có $3!$ cách xếp.
			\item Xếp nhóm $\mathrm{N}$ cùng với $9$ nam, có  $10!$ cách xếp.
		\end{itemize}
		Suy ra có $3! \cdot  10!$ cách xếp.\\
	Vậy xác suất cần tìm là  $\dfrac{3! \cdot 10!}{12!}=\dfrac{1}{22}$.}
\end{ex}

% nguồn: MaHoaNH_DungSai_Lop10
\begin{ex}%[0D0V2-3]%[ma:bc3f149c]
	Xếp ngẫu nhiên $4$ học sinh gồm $2$ nam và $2$ nữ vào hai dãy ghế đối diện nhau,  mỗi dãy có $2$ ghế. Tính xác suất $P$ để $2$ học sinh nam cùng ngồi vào một dãy ghế.
	 \choiceTF
		{$P=\dfrac{1}{6}$} 
		{$P=\dfrac{1}{12}$} 
		 {$P=\dfrac{2}{3}$}
		 {\True $P=\dfrac{1}{3}$}
	\loigiai{ 
		Gọi $A$ là biến cố  \lq\lq $2$ học sinh nam cùng ngồi vào một dãy ghế \rq\rq. \\
		$\bullet$ Xác định số phần tử của không gian mẫu\\ 
		Xếp $4$ trong $4$ học sinh vào  dãy ghế vừa chọn có $4!=24$ cách.\\
		$\bullet$ Xác định số phần tử của biến  cố $A$\\
		Chọn $1$ dãy ghế trong $2$  dãy ghế có $2$ cách.\\
		 Xếp $2$ học sinh nam vào dãy ghế vừa chọn có $2$ cách.\\
		Xếp $2$ học  sinh nữ vào dãy ghế còn lại có $2$ cách.\\
		Vậy số phần tử  của biến cố $A$ là $2\cdot 2\cdot 2=8$.\\
		 Xác suất để $2$ hai học sinh nam cùng ngồi một dãy ghế là $P=\dfrac{8}{24}=\dfrac{1}{3}$.
	 }
\end{ex}

% nguồn: MaHoaNH_DungSai_Lop10
\begin{ex}%[0D0V2-3]%[ma:94796687]
	Xét phép thử: Xếp ngẫu nhiên $7$ học sinh $A$, $B$, $C$, $D$, $E$,  $F$, $G$ vào một chiếc bàn dài có đúng $7$ ghế.
	\choiceTF{\True  Số phần tử của không gian mẫu là $n(\Omega)=7!$}{
		Xác suất  để $D$ không ngồi đầu bàn là $\dfrac{6}{7}$
		}{\True Xác suất để $A$ và $B$ không ngồi cạnh nhau là $\dfrac{5}{7}$}{Xác suất để $A$ và $B$ ngồi cạnh nhau, đồng thời $C$, $D$ không ngồi cạnh nhau và không ngồi cạnh $A,  B$ là $\dfrac{4}{21}$}
	\loigiai{ 
		\begin{itemchoice}
			\itemch {\bf Đúng.} Số phần tử của không  gian mẫu là $n(\Omega)=7!$.
			\itemch {\bf Sai.} Gọi $X$ là biến cố \lq\lq$D$ không  ngồi đầu bàn\rq\rq.
			\begin{enumerate}
				\item Bước 1: Xếp $D$ ngồi vào ghế không phải  đầu bàn, có $5$ cách.
				\item Bước 2: Xếp các bạn còn lại, có  $6!$ cách.
			\end{enumerate}
			Suy ra  $n(X)=5\cdot 6!$ cách.\\
			Xác suất $\mathrm{P}(X)=\dfrac{5\cdot  6!}{7!}=\dfrac{5}{7}$.
			\itemch {\bf Đúng.}  Gọi $Y$ là biến cố \lq\lq $A$ và $B$ không ngồi cạnh nhau\rq\rq. \\
			Xếp $5$ bạn $C$, $D$, $E$, $F$, $G$ vào bàn dài  có $5!$ cách.\\
			Giữa $5$ bạn tạo thành $6$ khoảng trống, chọn ra $2$ khoảng  trống và xếp $2$ bạn $A$, $B$ vào đó, có $\mathrm{A}_6^2$ (cách).\\
			Vậy $n(Y)=5!\cdot  \mathrm{A}_6^2$, suy ra $\mathrm{P}(Y)=\dfrac{5}{7}$.
			\itemch {\bf Sai.} Gọi $Z$ là biến cố \lq\lq $A$ và $B$ ngồi cạnh nhau, đồng thời $C$, $D$ không ngồi cạnh nhau và không ngồi cạnh  $A, B$\rq\rq. \\
			Xếp $3$ bạn $E$, $F$, $G$  vào bàn dài có $3!$ cách.\\
			Giữa $3$ bạn tạo thành $4$ khoảng trống, chọn ra $3$ khoảng trống  để xếp $AB$, $C$ và $D$, có $\mathrm{A}_4^3$ (cách).\\
			 Đổi chỗ $AB$ cho nhau có $2!$ cách.\\
			Vậy $n(Z)=3!\cdot \mathrm{A}_4^3\cdot 2!$,  suy ra $\mathrm{P}(Z)=\dfrac{3!\cdot \mathrm{A}_4^3\cdot 2!}{7!}=\dfrac{2}{35}$.
		\end{itemchoice}
	 }
\end{ex}

% nguồn: MaHoaNH_DungSai_Lop10
\begin{ex}%[0D0V2-3]%[ma:f11d1b10]
	Chọn ngẫu nhiên $4$ bông từ $4$  bông đỏ, $5$ bông xanh và $6$ bông vàng.
	\choiceTF 
		{ Xác suất $4$ bông được chọn cùng màu là  $\dfrac{5}{91}$}
		{\True Xác  suất $4$ bông được chọn có $2$ bông đỏ, $1$ bông xanh, $1$ bông vàng là $\dfrac{12}{91}$}
		{ Xác suất $4$ bông được chọn có  ít nhất $1$ bông đỏ là $\dfrac{22}{91}$}
		{\True Xác suất $4$ bông được chọn có  đủ ba màu là $\dfrac{48}{91}$}
	\loigiai{ 
		\begin{itemchoice}
			\itemch \textbf{Sai.}\\ 
			Số phần tử của không gian mẫu: $n( \Omega)=\mathrm{C}_{15}^4=1\,365$.\\ 
			Gọi $A$ là biến cố \lq\lq Bốn bông được chọn cùng màu \rq\rq\  thì $n( A)=\mathrm{C}_{4}^4+\mathrm{C}_{5}^4+\mathrm{C}_{6}^4=21$.\\
			Xác suất của  biến cố $A$ là $\mathrm{P}( A)=\dfrac{n( A)}{n( \Omega)}=\dfrac{1}{65}$.
			\itemch  \textbf{Đúng.}\\
			Gọi $B$ là biến  cố \lq\lq Bốn bông được chọn có $2$ bông đỏ, $1$ bông xanh, $1$ bông vàng \rq\rq.\\
			Khi đó $n( B)=\mathrm{C}_{4}^2\cdot\mathrm C_{5}^{1}\cdot  \mathrm{C}_{6}^{1}=180$.\\
			Xác suất của biến cố  $B$ là $\mathrm{P}( B)=\dfrac{n( B)}{n( \Omega)}=\dfrac{12}{91}$.
			 \itemch \textbf{Sai.}\\
			Gọi $C$ là biến cố \lq\lq Bốn bông được chọn có ít nhất $1$ bông đỏ\rq\rq.  Khi đó $\overline{C}$ là biến cố \lq\lq Bốn bông được chọn không có bông đỏ \rq\rq.\\
			Suy ra $n(  \overline{C})=\mathrm{C}_{11}^4=330$.\\
			Vậy xác suất của biến cố $C$ là $\mathrm{P}( C)=1-\mathrm{P}(  \overline{C})=1-\dfrac{330}{1\,365}=\dfrac{69}{91}$.
			\itemch  \textbf{Đúng.}\\
			Gọi $D$ là biến cố \lq\lq Bốn bông được chọn có đủ ba  màu \rq\rq. Khi đó ta có các trường hợp sau:\\
			\textbf{Trường hợp 1}: Chọn được 2 bông đỏ,  1 bông xanh và 1 bông vàng có $\mathrm{C}_{4}^2 \cdot \mathrm{C}_{5}^{1}\cdot \mathrm{C}_{6}^{1} =180$ cách.\\
			\textbf{Trường hợp 2}: Chọn được 1 bông đỏ, 2 bông  xanh và 1 bông vàng có $\mathrm{C}_{4}^1 \cdot \mathrm{C}_{5}^{2}\cdot \mathrm{C}_{6}^{1} =240$ cách.\\
			\textbf{Trường hợp 3}: Chọn được 1  bông đỏ, 1 bông xanh và 2 bông vàng có $\mathrm{C}_{4}^1 \cdot \mathrm{C}_{5}^{1}\cdot \mathrm{C}_{6}^{2} =300$ cách.\\
			Theo  quy tắc cộng, ta có $n( D)=180+240+300=720$.\\
			Xác suất  của biến cố $D$ là $\mathrm{P}( D)=\dfrac{n( D)}{n( \Omega)}=\dfrac{48}{91}$.
		\end{itemchoice}
	 }
\end{ex}


%%%%%%%%%%%%%%%% TRẢ LỜI NGẮN %%%%%%%%%%%%%%%%

% nguồn: MaHoaNH_TraLoiNgan_Lop10
\begin{bt}%[0D0V2-3]%[ma:c0492c03]
	Mật khẩu mở máy tính của Dũng gồm $8$ kí tự, trong đó có $2$ kí tự đầu là chữ số, $6$ kí tự sau là các chữ cái thuộc tập hợp $\left\{ A,B,C,D,E \right\}$. Không may Dũng quên mất $3$ kí tự đầu tiên. Dũng chọn ra $2$  chữ số và một chữ cái thuộc tập hợp trên một cách ngẫu nhiên và thử mở máy tính. Xác suất để Dũng mở được máy tính có dạng $\dfrac{a}{b}$. Tính $a+b$.
	 \shortans{451}
	\loigiai{ 
		\textbf{Công đoạn $1$:} Chọn $2$ số  từ $10$ số và có phân biệt thứ tự nên có $\mathrm{A}_{10}^2$ cách.\\
		\textbf{Công đoạn $2$:} Chọn ra $1$ chữ từ $5$ chữ cái có $\mathrm{C}_5^1$  cách.\\
		Vậy theo quy tắc nhân có: $\mathrm{A}_{10}^2\cdot\mathrm{C}_5^1=450$ cách.\\ 
		 Vậy số phần tử không gian mẫu là $ n(\Omega)=450$.\\
		Gọi biến cố $A$ \lq\lq Xác  xuất để Dũng mở được máy\rq\rq.\\
		Do chỉ có $1$ mật  mã đúng duy nhất nên $n(A)=1$.\\
		 Vậy $P(A)=\dfrac{n(A)}{n(\Omega)}=\dfrac{1}{450}$.
	} 
\end{bt}

% nguồn: MaHoaNH_TraLoiNgan_Lop10
\begin{bt}%[0D0V2-3]%[ma:c5a2eee8]
	Một nhóm $10$ thanh niên gồm $6$ nam trong đó có  Quang, và $4$ nữ trong đó có Huyền được xếp ngẫu nhiên vào một dãy $10$ ghế để tham dự buổi phát động phong trào thanh niên tình nguyện. Xác suất để xếp được giữa $2$ bạn nữ gần nhau có đúng $2$ bạn nam, đồng thời Quang không ngồi cạnh Huyền là $\dfrac{1}{a}$, giá trị của $a$ là bao nhiêu?
	\shortans{280} 
	\loigiai{ 
		Giả sử  các ghế được đánh số từ $1$ đến $10$.\\
		Ta  có $n(\Omega)=10!$.\\
		Gọi $A$ là biến cố \lq\lq Giữa $2$ ban nữ gần nhau có  đúng $2$ bạn nam, đồng thời Quang không ngồi cạnh Huyền\rq\rq.\\
		Để có cách xếp  sao cho giữa $2$ bạn nữ có đúng $2$ bạn nam thì các bạn nữ phải ngồi ở các ghế đánh số $1$, $4$, $7$, $10$.\\
		Ta tính số cách sắp xếp chỗ ngồi sao cho  Huyền và Quang không ngồi cạnh nhau.
		\begin{itemize}
			\item Nếu Huyền ngồi ở ghế $1$ hoặc $10$ thì có $5$ cách xếp chỗ ngồi cho  Quang.
			\item Nếu Huyền ngồi ở ghế $4$ hoặc $7$ thì có $4$ cách  xếp chỗ ngồi cho Quang.
		\end{itemize}
		Do đó,  số cách xếp chỗ cho Quang và Huyền không ngồi cạnh nhau là $2\cdot 5+2\cdot 4=18$.\\
		Số cách xếp chỗ  ngồi cho $3$ nữ còn lại và $5$ nam còn lại là $3!\cdot 5!$.
		$$ 
		\Rightarrow n(A)=18\cdot 3  !\cdot 5 !=12960\Rightarrow P(A)=\dfrac{n(A)}{n(\Omega)}=\dfrac{12960}{10 !}=\dfrac{1}{280}.
		$$ 
		Vây xác suất  cần tìm là $\dfrac{1}{280}$.
	} 
\end{bt}


%%%%%%%%%%%%%%%% TỰ LUẬN %%%%%%%%%%%%%%%%

% nguồn: Ok_TuanLuan_Toan11_Lan1
\begin{bt}%[0D0Y2-3]%[ma:e9dc4e8b]
	Khảo sát một trường trung học phổ thông, người ta thấy có $20\%$ học sinh thuận tay trái và $35\%$ học sinh bị cận thị. Giả sử đặc điểm thuận tay nào không  ảnh hưởng đến việc học sinh có bị cận thị hay không. Gặp ngẫu nhiên một học sinh của trường. Tính xác suất của biến cố học sinh đó bị cận thị hoặc thuận tay trái.
	\loigiai{Gọi $A$ là biến cố học sinh thuận tay  trái, khi đó $P(A)=0{,}2$.\\
		Gọi $B$ là biến cố học sinh cận thị, khi đó $P(B)=0{,}35$.\\ 
	Khi đó xác suất của biến cố  học sinh đó bị cận thị hoặc thuận tay trái là $P(A\cup B)=P(A)+P(B)-P(AB)=0{,}48$.}
\end{bt}

% nguồn: Ok_TuanLuan_Toan11_Lan1
\begin{bt}%[0D0Y2-3]%[ma:f842326e]
	Tỉ  lệ nảy mầm của một loại hạt giống là $0{,}8$. Gieo hai hạt giống một cách độc lập nhau. Tính xác suất để có đúng một trong hai hạt giống đó nảy mầm.
	\loigiai{ 
	Xác suất để hạt giống nảy mầm là $0{,}8$ và xác suất để hạt giống  không nảy mầm là $0{,}2$, khi đó xác suất để có đúng một trong hai hạt giống đó nảy mầm là $0{,}8\times 0{,}2+0{,}2\times 0{,}8=0{,}32$.}
\end{bt}

% nguồn: Ok_TuanLuan_Toan11_Lan1
\begin{bt}%[0D0K2-3]%[ma:98c8534c]
	Một hộp chứa  $5$ quả bóng xanh, $6$ quả bóng đỏ và $2$ quả bóng vàng có cùng kích thước và khối lượng. Chọn ra ngẫu nhiên từ hộp $3$ quả bóng. Tính xác suất của các biến cố:
	\begin{enumerate}[a)]
		\item "Cả ba quả  bóng lấy ra đều có cùng màu"
		\item "Có ít nhất $2$ quả bóng xanh trong $3$ quả bóng  lấy ra"
	\end{enumerate}
	 \loigiai{
		\begin{enumerate}[a)]
			\item "Cả ba quả bóng lấy ra đều  có cùng màu"\\
			\textbf{TH1:}  Xác suất để lấy được ba quả bóng màu xanh là $P_1=\dfrac{\mathrm{C}_5^3}{\mathrm{C}_{13}^3}$.\\
			\textbf{TH2:} Xác suất để lấy được ba quả bóng  màu đỏ là $P_2=\dfrac{\mathrm{C}_6^3}{\mathrm{C}_{13}^3}$.\\
			Vậy  xác suất để lấy được ba quả bóng đều có cùng màu là $P=P_1+P_2=\dfrac{15}{143}$.
			 \item "Có ít nhất $2$ quả bóng xanh trong $3$ quả bóng lấy ra"\\
			\textbf{TH1:} Xác suất để  lấy được ba quả bóng màu xanh là $P_1=\dfrac{\mathrm{C}_5^3}{\mathrm{C}_{13}^3}$.\\
			\textbf{TH2:} Xác suất để lấy được hai quả bóng màu  xanh, và một quả màu đỏ hoặc 1 quả màu vàng là $P_2=\dfrac{\mathrm{C}_5^2\times(6+2)}{\mathrm{C}_{13}^3}$.\\
			Vậy xác suất để lấy được Có  ít nhất $2$ quả bóng xanh trong $3$ quả bóng là $P=P_1+P_2=\dfrac{45}{143}$.
		\end{enumerate}
	 }
\end{bt}

% nguồn: Ok_TuanLuan_Toan11_Lan2
\begin{ex}%[0D0Y2-3]%[ma:d20bf300]
	 Khảo sát một trường trung học phổ thông, người ta thấy có $20\%$ học sinh thuận tay trái và $35\%$ học sinh bị cận thị. Giả sử đặc điểm thuận tay nào không ảnh hưởng đến việc học sinh có bị cận thị hay không. Gặp ngẫu nhiên một học sinh của trường. Tính xác suất của biến cố học sinh đó bị cận thị hoặc thuận tay trái.
	\loigiai{Gọi $A$ là biến cố học sinh thuận  tay trái, khi đó $P(A)=0{,}2$.\\
		Gọi $B$ là biến cố học sinh cận thị, khi đó $P(B)=0{,}35$.\\ 
	Khi đó  xác suất của biến cố học sinh đó bị cận thị hoặc thuận tay trái là $P(A\cup B)=P(A)+P(B)-P(AB)=0{,}48$.}
\end{ex}

% nguồn: Ok_TuanLuan_Toan11_Lan2
\begin{ex}%[0D0Y2-3]%[ma:3f484a69]
	Tỉ lệ nảy mầm của một loại hạt giống là $0{,}8$. Gieo hai hạt giống một cách độc lập  nhau. Tính xác suất để có đúng một trong hai hạt giống đó nảy mầm.
	\loigiai{ 
	Xác suất để hạt giống nảy mầm là $0{,}8$ và xác suất để hạt giống không nảy mầm là $0{,}2$, khi đó xác suất để có đúng một trong hai hạt giống đó nảy mầm là $0{,}8\times  0{,}2+0{,}2\times 0{,}8=0{,}32$.}
\end{ex}

% nguồn: Ok_TuanLuan_Toan11_Lan2
\begin{ex}%[0D0K2-3]%[ma:58169b17]
	Một hộp chứa $5$ quả bóng xanh, $6$ quả bóng đỏ và $2$ quả bóng vàng có cùng kích thước và khối lượng. Chọn ra ngẫu nhiên từ hộp $3$ quả bóng.  Tính xác suất của các biến cố:
	\begin{enumerate}[a)]
		\item "Cả ba quả bóng lấy  ra đều có cùng màu"
		\item "Có ít  nhất $2$ quả bóng xanh trong $3$ quả bóng lấy ra"
	\end{enumerate}
	\loigiai{ 
		\begin{enumerate}[a)]
			\item "Cả ba quả bóng lấy ra đều có cùng  màu"\\
			\textbf{TH1:} Xác suất để lấy được  ba quả bóng màu xanh là $P_1=\dfrac{\mathrm{C}_5^3}{\mathrm{C}_{13}^3}$.\\
			 \textbf{TH2:} Xác suất để lấy được ba quả bóng màu đỏ là $P_2=\dfrac{\mathrm{C}_6^3}{\mathrm{C}_{13}^3}$.\\
			Vậy xác suất để lấy được  ba quả bóng đều có cùng màu là $P=P_1+P_2=\dfrac{15}{143}$.
			\item "Có ít nhất $2$ quả bóng xanh trong $3$ quả  bóng lấy ra"\\
			\textbf{TH1:} Xác suất để lấy  được ba quả bóng màu xanh là $P_1=\dfrac{\mathrm{C}_5^3}{\mathrm{C}_{13}^3}$.\\
			\textbf{TH2:} Xác suất để lấy được hai quả bóng màu xanh, và một quả màu đỏ hoặc 1 quả màu vàng là $P_2=\dfrac{\mathrm{C}_5^2\times(6+2)}{\mathrm{C}_{13}^3}$.\\ 
			Vậy xác suất để lấy được Có ít  nhất $2$ quả bóng xanh trong $3$ quả bóng là $P=P_1+P_2=\dfrac{45}{143}$.
		\end{enumerate}
	} 
\end{ex}
